\scanpage{80}
% archive sheet number: 79
% archive/scan mark: right-angle registration-like mark at edge (upright upper left)
% Typescript upside down in scan; transcribed in its upright reading order.
\uline{STATUTS de l'UNIVERSITE DES SCIENCES ET TECHNIQUES DU LANGUEDOC}

\uline{Article 20} (Nouvelle rédaction approuvée par le conseil de l'Université le 29 juin 1977).

\begin{center}
--

\fbox{CONSEIL SCIENTIFIQUE}
\end{center}

I.- \uline{MISSION}.

Conformément aux dispositions de l'article 13 de la loi d'orientation de l'enseignement supérieur, un conseil scientifique détermine les programmes de recherche et répartit les crédits correspondants. Pour cela :
\begin{itemize}[label=-]
\item il garantit le libre arbitre des universitaires dans le choix des axes de recherche,
\item il favorise la participation des équipes de l'Université à des travaux s'insérant dans les axes des programmes de recherche internationaux, nationaux et régionaux.
\item il établit les bilans et propose les choix en ce qui concerne la politique de recherche scientifique, compte tenu du potentiel de recherche des laboratoires, des axes de recherche de l'Université et de la politique définie par les organismes internationaux, nationaux et régionaux.
\item il répartit les crédits de fonctionnement de recherche non affectés de la façon suivante :
\begin{itemize}[label=$\cdot$]
\item une partie attribuée en tenant compte du nombre et de l'encadrement des laboratoires;
\item une partie attribuée sur programme. La partie attribuée sur programme est utilisée notamment pour :
\begin{itemize}[label=-]
\item aider les jeunes équipes pour leur permettre d'atteindre le seuil de prise en charge par les organismes nationaux;
\item assurer les équipements d'intérêt collectif (laboratoire ou service commun);
\item apporter son soutien aux formations ayant un programme reconnu sur le plan national, ou dont il reconnaîtra l'intérêt;
\item mettre en place des opérations susceptibles d'être soutenues et relayées par les organismes nationaux et régionaux (C.N.R.S., D.G.R.S.T., etc.).
\end{itemize}
\end{itemize}
\end{itemize}
